% Generated draft: requires mathematical and source review.
\documentclass[11pt]{article}
\usepackage[margin=25mm]{geometry}
\usepackage{fontspec}
\setmainfont{Latin Modern Roman}
\usepackage{amsmath,amssymb,mathtools,unicode-math}
\setmathfont{Latin Modern Math}
\usepackage{enumitem,needspace,xurl,hyperref}
\setlength{\parindent}{0pt}
\setlength{\parskip}{5pt}
\setlength{\emergencystretch}{3em}
\newcommand{\N}{\mathbb{N}}
\newcommand{\Z}{\mathbb{Z}}
\newcommand{\Q}{\mathbb{Q}}
\newcommand{\R}{\mathbb{R}}
\newcommand{\C}{\mathbb{C}}
\newcommand{\F}{\mathbb{F}}
\newcommand{\A}{\mathbb{A}}
\newcommand{\PP}{\mathbb{P}}
\newcommand{\E}{\mathbb{E}}
\newcommand{\eps}{\varepsilon}
\newcommand{\dd}{\,\mathrm{d}}
\begin{document}
\section*{UnsolvedMath ID: 20003349}
\textbf{Source problem code:} catalogue entry
\subsection{Definitions and mathematical statement}
\textbf{Question 50 (B. Weiss).} Let \(X=G/\Gamma\), let \(A\) be the torus appearing in the missing preceding context, with \(\dim A=r\ge2\), and take \(x\in X\). Suppose that for every nontrivial one-parameter subgroup \(D\le A\), the orbit \(Dx\) is not divergent. Does there exist a compact subset \(K\subset X\) such that\[\limsup_{T\to\infty}\frac{\lambda(\{a\in A_T:ax\in K\})}{\lambda(A_T)}>0?\] Here \(\lambda\) is Haar measure on \(A\). The excerpt omits the preceding definition of \(A_T\) and the precise ambient-group hypotheses. The aid uses the norm-ball interpretation: identify the connected real split torus \(A\) with \((\R^r,+)\) by its logarithm, fix a norm, and set \(A_T=\{\exp(v):\|v\|\le T\}\). Divergence of a one-parameter orbit means properness of its orbit map, equivalently eventual departure from every compact set as the parameter leaves compact subsets of \(\R\). Nondivergence is not the assertion that the whole orbit is relatively compact. These conventions are explicit reconstructions, pending the missing source context.\par The attached aid supplies three boundary statements, rather than a solution in the general homogeneous setting. For a continuous \(A\cong\R^r\) action on a locally compact second-countable Hausdorff space, put\[\eta_T=\frac1{\lambda(A_T)}\int_{A_T}\delta_{ax}\,d\lambda(a).\] The existence of a compact \(K\) with positive upper return proportion is equivalent to the existence of a nonzero vague subsequential limit of \(\eta_T\), where vague convergence tests \(C_c(X)\). Every such limit is an \(A\)-invariant finite Radon measure with total mass at most one. Nonzero mass does not assert mass one or full equidistribution. In the homogeneous case with \(\Gamma\) discrete, the stabilizer \(H=\{a\in A:ax=x\}\) is discrete. If \(Ax\) is closed, so that the orbit is identified with \(A/H\), nondivergence of every one-parameter subgroup forces \(H\) to have rank \(r\); hence \(Ax\) is compact and choosing \(K=Ax\) gives return proportion one. Finally, the aid reports a nonhomogeneous translation-hull example of a locally compact \(\R^2\)-action and a point for which every line returns at unbounded parameter values to one common compact target, although every compact target has norm-ball return proportion tending to zero. The actual hull construction is not supplied; this example is a reported obstruction to inferring density from purely directional recurrence in arbitrary actions, not a homogeneous counterexample to Weiss's question.\par \textbf{Boundary of the extracted record.} The text after Question 50 begins Section 9, \emph{Andr\'e--Oort Conjecture}, and is separate background. It describes a Shimura datum \((G,\mathcal X)\) with \(G/\Q\) reductive and \(\mathcal X\) a \(G(\R)\)-conjugacy class of homomorphisms \(h:\C^\times\to G(\R)\): the adjoint characters on the complexified Lie algebra are \(z/\bar z,1,\bar z/z\), conjugation by \(h(i)\) is a Cartan involution, and the extracted third condition excludes factors with trivial adjoint action. Here \(G^{\mathrm{ad}}\) is the quotient by the center. With finite adeles \(\A_f\) and compact open \(K\le G(\A_f)\), it defines\[\operatorname{Sh}_K(G,\mathcal X)=G(\Q)\backslash(\mathcal X\times G(\A_f))/K,\] recalls its complex algebraic structure through Baily--Borel and Hermitian locally symmetric components, and forms the projective limit over \(K\). A special point has \(h(\C^\times)\subset T(\R)\) for a torus \(T/\Q\). Its example is \(G=\mathrm{GL}_2\), \(h(a+bi)=\left(\begin{smallmatrix}a&b\\-b&a\end{smallmatrix}\right)\), \(K=\mathrm{GL}_2(\widehat\Z)\), giving the modular curve and complex-multiplication elliptic curves as special points. Morphisms of Shimura data induce maps of these spaces. The projections from level \(K\cap gKg^{-1}\) to levels \(K\) and \(gKg^{-1}\), followed by the right-\(g\) identification with level \(K\), define the finite-fiber Hecke correspondence \(T_g\). The source calls irreducible components of Hecke translates of images of Shimura subdata special subvarieties, or subvarieties of Hodge type, and reports density of their special points in the Zariski and analytic topologies. Its final sentence introduces an ensuing converse conjecture but does not state it. None of this changes the quantifiers in Question 50; the extracted Shimura-datum conventions and quotient terminology require their own source review.
\subsection{Short English statement}
The question asks whether nondivergence in every one-dimensional direction forces a positive fraction of a large torus ball to return to some compact set. The aid establishes the closed-orbit case and characterizes the desired conclusion by nonzero invariant vague limits. Its general-action obstruction does not answer the homogeneous question. Appended Shimura-variety background belongs to the next section.
\subsection{Sources}
\begin{itemize}\item Local UnsolvedMath registry, permanent ID 20003349, code catalogue entry; canonical \texttt{aim-other-notes.json}, \texttt{notes[155]}, question tag.\item American Institute of Mathematics Problem Lists, Other domain, \emph{Emerging applications of measure rigidity}, Question 50 attributed to B. Weiss. \url{https://aimath.org/WWN/measrigid/measrigid.pdf}. The extraction continues into Section 9, \emph{Andr\'e--Oort Conjecture}, including a Baily--Borel attribution but no complete bibliographic entry.\item Attached research attempt 1, \emph{Directional nondivergence, invariant cluster measures, and the closed-orbit case}, recorded 2026-08-14: \texttt{partial\_result}, problem-at-run \texttt{open}, rigor \texttt{proved}, novelty confidence low; registry \texttt{partially\_solved}, classification \texttt{PARTIAL-PROGRESS}. Statement status \texttt{reconstructed\_unverified}. The machine-generated aid reports inspection of the official PDF, but that inspection is not independently verified here; the aid is not peer reviewed and must be checked before citation.\end{itemize}
\end{document}
